% ============================================================
% Bản 45 phút -- Slide 29-32
% Nguồn: partADC_7.tex (Slide 29) + part5.tex (Slide 30-32)
% CHỈ TẠO MỚI trong Slides45min/, không sửa file gốc.
% ============================================================

% --- Slide 29 (gốc: partADC_7.tex, "Cải tiến 5: final prune trên toy — khép lại chương ADC") ---
% Đã bỏ minh hoạ "final prune trên toy" + kết luận N nhỏ hơn: cả prune multinomial
% (pruning_score=None) và final_prune_structgs (lời gọi bị comment, train.py:419)
% KHÔNG chạy trong pipeline thật, nên không được trình bày như kết quả huấn luyện.
\begin{frame}{Tổng kết cải tiến ADC của SADGS}
\footnotesize
\begin{itemize}
  \item \textbf{Đang chạy trong pipeline thật}: (1) Importance đa view AND với gradient; (2) split theo gradient tuyệt đối; (3) trần opacity $\alpha\leftarrow\min(\alpha,0.8)$ sau mỗi vòng densify+prune.
  \item \textbf{Không chạy trong pipeline thật} (có trong code nhưng vô hiệu): prune theo lấy mẫu multinomial -- \texttt{train.py} luôn truyền \texttt{pruning\_score=None}; và \texttt{final\_prune\_structgs} -- lời gọi bị \textbf{comment out} tại \texttt{train.py:419}. Tại \texttt{prune\_iterations}, code thật chỉ xoá theo $\alpha<0.1$.
\end{itemize}
\end{frame}

% --- Slide 30 (gốc: part5.tex, "Vì sao 3DGS gốc chậm?") ---
\begin{frame}{Vì sao 3DGS gốc chậm?}
\begin{itemize}
    \item Mỗi iteration gồm: \textbf{render} (rasterize), \textbf{tính loss}, \textbf{backward}, định kỳ \textbf{densify/prune}.
    \item Với hàng trăm nghìn -- hàng triệu Gaussian, \textbf{rasterize} và \textbf{backward} chiếm phần lớn thời gian mỗi bước.
\end{itemize}
\vspace{0.3cm}
\centering
\footnotesize
\begin{tabular}{lcc}
\toprule
\textbf{Khâu} & \textbf{Tần suất} & \textbf{Tỉ trọng thời gian} \\
\midrule
Render (project + rasterize) & mỗi iteration & \textasciitilde 45\% \\
Backward (gradient) & mỗi iteration & \textasciitilde 35\% \\
Tính loss & mỗi iteration & \textasciitilde 10\% \\
Densify/Prune & định kỳ & \textasciitilde 10\% (đột biến) \\
\bottomrule
\end{tabular}
\end{frame}

% --- Slide 31 (gốc: part5.tex, "Ba đòn bẩy tăng tốc của SADGS (tổng quan)") ---
\begin{frame}{Ba đòn bẩy tăng tốc của SADGS (tổng quan)}
\begin{columns}[t]
\column{0.33\linewidth}
\centering
\textbf{1. Multi-view score}
\begin{itemize}
    \footnotesize
    \item Đánh giá Gaussian qua nhiều góc nhìn
    \item Pruning thông minh hơn
\end{itemize}
\column{0.33\linewidth}
\centering
\textbf{2. Compact Box}
\begin{itemize}
    \footnotesize
    \item Giảm hệ số nhân bounding-box
    \item Giảm số tile chạm
\end{itemize}
\column{0.33\linewidth}
\centering
\textbf{3. Sparse Adam}
\begin{itemize}
    \footnotesize
    \item Lịch cập nhật phân tầng
    \item Giảm tần suất optimizer.step()
\end{itemize}
\end{columns}
\vspace{0.4cm}
\begin{itemize}
    \item Cả ba đòn bẩy nhắm đúng khâu chiếm tỉ trọng lớn nhất: rasterize, densify/prune, cập nhật tham số.
\end{itemize}
\end{frame}

% --- Slide 32 (gốc: part5.tex, "Hiệu năng tích luỹ: Biểu đồ thác nước") ---
\begin{frame}{Hiệu năng tích luỹ: Biểu đồ thác nước}
\footnotesize
\begin{itemize}
    \item Waterfall chart minh hoạ tăng tốc \textbf{tích luỹ} khi bật lần lượt từng đòn bẩy lên baseline 3DGS gốc.
    \item Thứ tự: baseline $\to$ + Compact Box $\to$ + Multi-view score $\to$ + Sparse Adam $\to$ SADGS đầy đủ.
\end{itemize}
\begin{center}
\includegraphics[width=0.68\linewidth]{ch8_levers_waterfall.png}
\end{center}
\end{frame}
